%======================================================================
\section*{Chapter 9: Tests of Hypotheses Based on a Single Sample}
\addcontentsline{toc}{section}{Chapter 9: Tests of Hypotheses Based on a Single Sample}

\subsection*{9.1 Hypotheses and Test Procedures}
\addcontentsline{toc}{subsection}{9.1 Hypotheses and Test Procedures}

Example: \textbf{Null hypothesis} $H_0:\mu = 50$, \textbf{alternative hypothesis} $H_a:\mu \neq 50$.

There are two possible outcomes:
\begin{enumerate}
  \item If the evidence obtained from the sample strongly supports the alternative
        hypothesis $H_a$, then we \emph{reject} $H_0$.
  \item If the evidence from the sample does not support $H_a$, then we
        \emph{fail to reject} $H_0$.
\end{enumerate}

\textbf{Test procedure.} A test procedure has two components:
\begin{enumerate}
  \item A \textbf{test statistic}: a function of the sample. The decision is based on the test statistic.
  \item A \textbf{rejection region}: the set of computed values of the test statistic for which $H_0$ is rejected.
\end{enumerate}
The null hypothesis is rejected if the observed (computed) value of the test
statistic falls in the rejection region.

\textbf{The errors.} There are two potential errors in a statistical test procedure:
\begin{enumerate}
  \item \textbf{Type I error}: rejecting $H_0$ when it is true.
  \item \textbf{Type II error}: not rejecting $H_0$ when it is false.
\end{enumerate}
Their probabilities are denoted by $\alpha$ and $\beta$:
\begin{align*}
\alpha &= P(\text{type I error}) = P(H_0 \text{ is rejected when it is true}),\\
\beta  &= P(\text{type II error}) = P(H_0 \text{ is not rejected when it is false}).
\end{align*}
\begin{itemize}
  \item The probability of rejecting $H_0$ when it is true, $\alpha$, is called the \textbf{significance level}.
  \item The probability of rejecting $H_0$ when it is false is called the \textbf{power} of the test procedure; it equals $1-\beta$.
\end{itemize}

Suppose that the parameter of interest is $\theta$ and the null hypothesis is
$H_0:\theta=\theta_0$, where $\theta_0$ is the \emph{null value}. The alternative
hypothesis appears in one of three forms (the equality always belongs to $H_0$):
\begin{enumerate}
  \item $H_a:\theta > \theta_0$ \quad (implicitly $H_0:\theta \le \theta_0$)
  \item $H_a:\theta < \theta_0$ \quad (implicitly $H_0:\theta \ge \theta_0$)
  \item $H_a:\theta \neq \theta_0$ \quad ($H_0:\theta = \theta_0$)
\end{enumerate}

\textbf{Formulating hypotheses (example).} For a disease, the current drug gives an
average life time of 12 months. A new drug (Drug 2) is claimed to give an average
life time of 24 months. We test
\[
H_0:\mu = 12 \qquad\text{vs.}\qquad H_a:\mu > 12,
\]
where $\mu$ is the mean life time under the new drug, using e.g.\ a sample of 50 patients and their sample mean $\bar{x}$.

\subsection*{9.2 Tests About a Population Mean}
\addcontentsline{toc}{subsection}{9.2 Tests About a Population Mean}

\textbf{Case I: normal population, $\sigma$ known.}

Assumptions: (1) the population distribution is normal; (2) $\sigma$ is known.

Consider the standardized $\bar{X}$ under the assumption that $H_0$ is true:
\[
Z = \frac{\bar{X}-\mu_0}{\sigma/\sqrt{n}} \quad\longrightarrow\quad z = \frac{\bar{x}-\mu_0}{\sigma/\sqrt{n}}.
\]
Here $z$ measures the distance between $\bar{x}$ and $\mu_0$ in ``standard deviation units''.
Take $H_a:\mu>\mu_0$ (i.e.\ $H_0:\mu\le\mu_0$):
\begin{enumerate}
  \item \emph{Large positive $z$}: the difference between $\bar{x}$ and $\mu_0$ is large, so $z$ is positive and large. This is evidence for $H_a$.
  \item \emph{Small positive $z$}: $\bar{x}$ is only a little bit greater than $\mu_0$; we attribute this to random luck and fail to reject $H_0$ (the evidence is not strong enough).
  \item \emph{Negative $z$}: $\bar{x}<\mu_0$, we fail to reject $H_0$.
\end{enumerate}
We only reject $H_0$ if $\bar{x}$ is \emph{significantly} greater than $\mu_0$, i.e.\ $z$ must be
greater than or equal to some positive value $c$. The rejection region is $z \ge c$.

The cutoff value $c$ is used to control the probability of a type I error at the desired level $\alpha$:
\begin{align*}
\alpha &= P(H_0 \text{ rejected when true}) = P(Z\ge c \text{ when } Z\sim N(0,1)) = 1-\Phi(c)\\
&\Rightarrow\; \Phi(c)=1-\alpha \;\Rightarrow\; c=\Phi^{-1}(1-\alpha)=z_\alpha .
\end{align*}
So the rejection region is $z\ge z_\alpha$ (right tail, area $\alpha$). This test procedure is called \textbf{upper-tailed}.
\fig{rr_upper}{0.45}{Upper-tailed test: rejection region $z\ge z_\alpha$.}


Similarly, $H_a:\mu<\mu_0$ suggests the rejection region $z\le c=-z_\alpha$ (left tail, area $\alpha$);
this test is called \textbf{lower-tailed}.
\fig{rr_lower}{0.45}{Lower-tailed test: rejection region $z\le -z_\alpha$.}


Similarly, $H_a:\mu\ne\mu_0$ suggests the rejection region $z\le -c$ or $z\ge c$, i.e.\
either $z\le -z_{\alpha/2}$ or $z\ge z_{\alpha/2}$ (area $\alpha/2$ in each tail);
this test is called \textbf{two-tailed}.
\fig{rr_two}{0.45}{Two-tailed test: rejection region $|z|\ge z_{\alpha/2}$.}


\textbf{Example.} Suppose $\alpha=.05$ is desired for the two-tailed test. Then
\[
.05 = P(Z\le -c \text{ or } Z\ge c) = \Phi(-c) + 1-\Phi(c) = 2[1-\Phi(c)]
\;\Rightarrow\; c = z_{.025} = 1.96 .
\]

\textbf{The one-sample $z$ test.}
Null hypothesis $H_0:\mu=\mu_0$. Test statistic value: $z = \dfrac{\bar{x}-\mu_0}{\sigma/\sqrt{n}}$.

\begin{center}
\begin{tabular}{ll}
\toprule
Alternative hypothesis & Rejection region for level $\alpha$ test\\
\midrule
$H_a:\mu>\mu_0$   & $z\ge z_\alpha$ (upper-tailed)\\
$H_a:\mu<\mu_0$   & $z\le -z_\alpha$ (lower-tailed)\\
$H_a:\mu\ne\mu_0$ & either $z\le -z_{\alpha/2}$ or $z\ge z_{\alpha/2}$ (two-tailed)\\
\bottomrule
\end{tabular}
\end{center}

\textbf{Power, $\beta$, and sample size determination for the one-sample $z$ test.}
Consider the upper-tailed test with rejection region $z\ge z_\alpha$. This is equivalent to
$\bar{x}\ge \mu_0 + z_\alpha\sigma/\sqrt{n}$, so $H_0$ is not rejected when
$\bar{x}<\mu_0+z_\alpha\sigma/\sqrt{n}$.
Let $\mu'$ be a particular value of $\mu$ that exceeds $\mu_0$. Then
\begin{align*}
\beta(\mu') &= P(H_0 \text{ is not rejected when it is false})
 = P\!\left(\bar{X}<\mu_0+z_\alpha\sigma/\sqrt{n} \text{ when } \mu=\mu'\right)\\
&= P\!\left(\frac{\bar{X}-\mu'}{\sigma/\sqrt{n}} < z_\alpha + \frac{\mu_0-\mu'}{\sigma/\sqrt{n}} \text{ when } \mu=\mu'\right)
 = \Phi\!\left(z_\alpha + \frac{\mu_0-\mu'}{\sigma/\sqrt{n}}\right).
\end{align*}
This is the probability of a type II error for a level $\alpha$ test.

To determine the sample size we need two restrictions: (1) $\beta(\mu')=\beta$ and
(2) $\alpha = P(\text{type I error})$. Then
\[
\Phi\!\left(z_\alpha + \frac{\mu_0-\mu'}{\sigma/\sqrt{n}}\right)=\beta
\;\Rightarrow\; z_\alpha + \frac{\mu_0-\mu'}{\sigma/\sqrt{n}} = \Phi^{-1}(\beta) = -z_\beta
\]
\[
\Rightarrow\; \frac{\mu_0-\mu'}{\sigma/\sqrt{n}} = -(z_\alpha+z_\beta)
\;\Rightarrow\; \sqrt{n} = \frac{\sigma(z_\alpha+z_\beta)}{\mu'-\mu_0}
\;\Rightarrow\; n = \left[\frac{\sigma(z_\alpha+z_\beta)}{\mu_0-\mu'}\right]^2
\]
for the upper-tailed procedure. In general:

\begin{center}
\begin{tabular}{ll}
\toprule
Alternative hypothesis & Type II error probability $\beta(\mu')$ for a level $\alpha$ test\\
\midrule
$H_a:\mu>\mu_0$   & $\Phi\!\left(z_\alpha + \dfrac{\mu_0-\mu'}{\sigma/\sqrt{n}}\right)$\\[2ex]
$H_a:\mu<\mu_0$   & $1-\Phi\!\left(-z_\alpha + \dfrac{\mu_0-\mu'}{\sigma/\sqrt{n}}\right)$\\[2ex]
$H_a:\mu\ne\mu_0$ & $\Phi\!\left(z_{\alpha/2} + \dfrac{\mu_0-\mu'}{\sigma/\sqrt{n}}\right) - \Phi\!\left(-z_{\alpha/2} + \dfrac{\mu_0-\mu'}{\sigma/\sqrt{n}}\right)$\\
\bottomrule
\end{tabular}
\end{center}
In each case the power is $1-\beta(\mu')$. The sample size is
\[
n = \begin{cases}
\left[\dfrac{\sigma(z_\alpha+z_\beta)}{\mu_0-\mu'}\right]^2 & \text{for a one-tailed test},\\[3ex]
\left[\dfrac{\sigma(z_{\alpha/2}+z_\beta)}{\mu_0-\mu'}\right]^2 & \text{for a two-tailed test (approximately)}.
\end{cases}
\]

\textbf{Case II: the one-sample $t$ test.}

Assumptions: (1) the population distribution is normal; (2) $\sigma$ is unknown.
Consider the random variable below under the assumption that $H_0$ is true:
\[
T = \frac{\bar{X}-\mu_0}{S/\sqrt{n}}, \qquad T\sim t_{n-1}.
\]

\textbf{Example.} Consider $H_0:\mu=\mu_0$ against $H_a:\mu>\mu_0$. Using the upper-tailed procedure,
\[
P(\text{type I error}) = P(H_0 \text{ rejected when true}) = P(T\ge t_{\alpha,n-1} \text{ when } T\sim t_{n-1}) = \alpha .
\]

\textbf{The one-sample $t$ test.}
Null hypothesis $H_0:\mu=\mu_0$. Test statistic value: $t = \dfrac{\bar{x}-\mu_0}{s/\sqrt{n}}$.

\begin{center}
\begin{tabular}{ll}
\toprule
Alternative hypothesis & Rejection region for level $\alpha$ test\\
\midrule
$H_a:\mu>\mu_0$   & $t\ge t_{\alpha,n-1}$ (upper-tailed)\\
$H_a:\mu<\mu_0$   & $t\le -t_{\alpha,n-1}$ (lower-tailed)\\
$H_a:\mu\ne\mu_0$ & either $t\le -t_{\alpha/2,n-1}$ or $t\ge t_{\alpha/2,n-1}$ (two-tailed)\\
\bottomrule
\end{tabular}
\end{center}

\textbf{Some practical advice.}
\begin{enumerate}
  \item For large $n$, the $t$ test is (approximately) accurate regardless of the underlying population distribution.
  \item For small $n$, it is not.
  \item Since the $z$ and $t$ distributions are quite similar for large $n$, the cutoffs obtained from the $z$ test will be similar to the cutoffs obtained from the $t$ test.
\end{enumerate}

\textbf{Power, $\beta$, and sample size determination for the one-sample $t$ test.}

\emph{Sample size.} For large $n$, the sample size formulas obtained for the $z$ test can be used.
Consider the one-tailed formula
\[
n = \left[\frac{\sigma(z_\alpha+z_\beta)}{\mu_0-\mu'}\right]^2 .
\]
To get rid of the unknown $\sigma$, we can
\begin{enumerate}
  \item replace $\sigma$ with $s$;
  \item use an estimate from a previous study;
  \item write $\mu'=\mu_0+d\sigma$, where $d$ is a positive or negative value.
\end{enumerate}
Substituting $\mu'=\mu_0+d\sigma$ we get
\[
n = \left[\frac{\sigma(z_\alpha+z_\beta)}{\mu_0-\mu_0-d\sigma}\right]^2 = \left[\frac{z_\alpha+z_\beta}{d}\right]^2 .
\]

\emph{Power and $\beta$.} For large $n$, we can use the $\beta(\mu')$ formulas from the $z$ test.
Consider $H_a:\mu>\mu_0$. Then
\[
\beta(\mu') = \Phi\!\left(z_\alpha + \frac{\mu_0-\mu'}{\sigma/\sqrt{n}}\right),
\]
and we can use solutions 1, 2, 3 above to get rid of the unknown $\sigma$. Following 3,
\[
\beta(\mu') = \beta(\mu_0+d\sigma) = \Phi\!\left(z_\alpha + \frac{\mu_0-\mu_0-d\sigma}{\sigma/\sqrt{n}}\right) = \Phi\!\left(z_\alpha - d\sqrt{n}\right).
\]

Recall that $\beta = P(\text{type II error}) = P(H_0 \text{ is not rejected when it is false})$.
When $H_0$ is false, $T=\dfrac{\bar{X}-\mu_0}{S/\sqrt{n}}$ does \emph{not} have a $t$ distribution; it has a more complicated distribution.

\textbf{Definition.} Let $Z\sim N(0,1)$ and $Y\sim\chi^2_\nu$ be independent random variables.
For any real number $\delta$, the random variable
\[
T = \frac{Z+\delta}{\sqrt{Y/\nu}}
\]
has a \textbf{noncentral $t$ distribution} with $\nu$ degrees of freedom and noncentrality parameter $\delta$.

If $\mu=\mu'$, then $T=\dfrac{\bar{X}-\mu_0}{S/\sqrt{n}}$ has a noncentral $t$ distribution with $n-1$ df and
$\delta = \dfrac{\mu'-\mu_0}{\sigma/\sqrt{n}}$.

\textbf{Example.} Determine the power of a lower-tailed test procedure. Assume that
$F(x;\nu,\delta)$ is the cdf of the noncentral $t$ distribution.

\emph{Solution.} The rejection region is $t\le -t_{\alpha,n-1}$, so
\begin{align*}
\text{Power} = 1-\beta(\mu') &= P\big(T\le -t_{\alpha,n-1} \text{ when } T \text{ has a noncentral } t \text{ distribution}\big)\\
&= P\!\left(T\le -t_{\alpha,n-1} \text{ when } T\sim \text{noncentral } t,\ \text{df}=n-1,\ \delta=\frac{\mu'-\mu_0}{\sigma/\sqrt{n}}\right)\\
&= F\!\left(-t_{\alpha,n-1};\ n-1,\ \frac{\mu'-\mu_0}{\sigma/\sqrt{n}}\right),
\end{align*}
which can be computed by software.
\fig{noncentral_t_power}{0.6}{Lower-tailed $t$ test: $\alpha$ is the area under the $t_{n-1}$ curve and the power is the area under the noncentral $t$ curve, both to the left of $-t_{\alpha,n-1}$.}


\subsection*{9.3 Tests About a Population Proportion}
\addcontentsline{toc}{subsection}{9.3 Tests About a Population Proportion}

\textbf{Large-sample test.} Consider $H_0:p=p_0$ when $n$ is large. If the null hypothesis is true, the random variable
\[
Z = \frac{\hat{p}-p_0}{\sqrt{p_0(1-p_0)/n}}
\]
has approximately a standard normal distribution.
Suppose $H_a:p>p_0$, which suggests the rejection region $z\ge z_\alpha$. Then
\[
P(\text{type I error}) = P(H_0 \text{ rejected when true}) = P(Z\ge z_\alpha \text{ when } Z\sim N(0,1) \text{ approx.}) \approx \alpha .
\]

\textbf{The one-proportion $z$ test.}
Null hypothesis $H_0:p=p_0$. Test statistic value: $z=\dfrac{\hat{p}-p_0}{\sqrt{p_0(1-p_0)/n}}$.

\begin{center}
\begin{tabular}{ll}
\toprule
Alternative hypothesis & Rejection region\\
\midrule
$H_a:p>p_0$   & $z\ge z_\alpha$ (upper-tailed)\\
$H_a:p<p_0$   & $z\le -z_\alpha$ (lower-tailed)\\
$H_a:p\ne p_0$ & either $z\le -z_{\alpha/2}$ or $z\ge z_{\alpha/2}$ (two-tailed)\\
\bottomrule
\end{tabular}
\end{center}
These procedures are valid when both $np_0\ge 10$ and $n(1-p_0)\ge 10$.

\textbf{Power, $\beta$, and sample size determination for the one-proportion $z$ test.}
\begin{itemize}
  \item When $H_0$ is true, $Z$ has approximately a standard normal distribution.
  \item When $H_0$ is false, $Z$ has approximately a normal distribution, with (where $p=p'$)
\end{itemize}
\[
E(Z) = E\!\left(\frac{\hat{p}-p_0}{\sqrt{p_0(1-p_0)/n}}\right) = \frac{E(\hat{p})-p_0}{\sqrt{p_0(1-p_0)/n}} = \frac{p'-p_0}{\sqrt{p_0(1-p_0)/n}},
\]
\[
V(Z) = V\!\left(\frac{\hat{p}-p_0}{\sqrt{p_0(1-p_0)/n}}\right) = \frac{V(\hat{p})}{p_0(1-p_0)/n} = \frac{p'(1-p')/n}{p_0(1-p_0)/n}.
\]
Write $q_0=1-p_0$ and $q'=1-p'$. Consider $H_a:p>p_0$ (rejection region $z\ge z_\alpha$). This is equivalent to the rejection region
$\hat{p}\ge p_0+z_\alpha\sqrt{p_0q_0/n}$. Then $H_0$ is not rejected when $\hat{p}<p_0+z_\alpha\sqrt{p_0q_0/n}$.
Now let $p'$ be a particular value of $p$ that exceeds $p_0$. Then
\begin{align*}
\beta(p') &= P(\text{type II error}) = P(H_0 \text{ is not rejected when } p=p')\\
&= P\!\left(\hat{p}<p_0+z_\alpha\sqrt{p_0q_0/n} \text{ when } p=p'\right)\\
&= P\!\left(\frac{\hat{p}-p'}{\sqrt{p'(1-p')/n}} < \frac{p_0-p'+z_\alpha\sqrt{p_0q_0/n}}{\sqrt{p'(1-p')/n}} \text{ when } p=p'\right)\\
&\approx \Phi\!\left(\frac{p_0-p'+z_\alpha\sqrt{p_0q_0/n}}{\sqrt{p'q'/n}}\right).
\end{align*}
The formulas for all three alternatives (textbook p.~528) are:

\begin{center}
\begin{tabular}{ll}
\toprule
Alternative hypothesis & $\beta(p')$\\
\midrule
$H_a:p>p_0$   & $\Phi\!\left(\dfrac{p_0-p'+z_\alpha\sqrt{p_0q_0/n}}{\sqrt{p'q'/n}}\right)$\\[3ex]
$H_a:p<p_0$   & $1-\Phi\!\left(\dfrac{p_0-p'-z_\alpha\sqrt{p_0q_0/n}}{\sqrt{p'q'/n}}\right)$\\[3ex]
$H_a:p\ne p_0$ & $\Phi\!\left(\dfrac{p_0-p'+z_{\alpha/2}\sqrt{p_0q_0/n}}{\sqrt{p'q'/n}}\right) - \Phi\!\left(\dfrac{p_0-p'-z_{\alpha/2}\sqrt{p_0q_0/n}}{\sqrt{p'q'/n}}\right)$\\
\bottomrule
\end{tabular}
\end{center}
For each $\beta(p')$ formula, the power is $1-\beta(p')$.

For the sample size we use two restrictions: (1) $\beta(p')=\beta$, (2) $\alpha\approx P(\text{type I error})$. Then
\[
\Phi\!\left(\frac{p_0-p'+z_\alpha\sqrt{p_0q_0/n}}{\sqrt{p'(1-p')/n}}\right)=\beta .
\]
Solving this equation for $n$ gives the sample size formulas (textbook p.~528):
\[
n = \begin{cases}
\left[\dfrac{z_\alpha\sqrt{p_0q_0}+z_\beta\sqrt{p'q'}}{p'-p_0}\right]^2 & \text{one-tailed test},\\[3ex]
\left[\dfrac{z_{\alpha/2}\sqrt{p_0q_0}+z_\beta\sqrt{p'q'}}{p'-p_0}\right]^2 & \text{two-tailed test (approximately)}.
\end{cases}
\]

\subsection*{9.4 $P$-Values}
\addcontentsline{toc}{subsection}{9.4 $P$-Values}

\textbf{Properties of the $P$-value.}
\begin{enumerate}
  \item The $P$-value is a probability.
  \item It is computed under the assumption that $H_0$ is true.
  \item It is a function of the sample data.
  \item To find the $P$-value, we identify all test statistic values that are as extreme as or more extreme than what we observed.
\end{enumerate}

\textbf{Example.} Suppose $\alpha=.01$ with the upper-tailed $z$ test.
\begin{itemize}
  \item Case I: $z\ge z_\alpha = z_{.01}=2.33$. Then the area to the right of $z$ is at most $.01$, i.e.\ $P$-value $\le .01$. Decision: reject $H_0$.
  \item Case II: $z< z_{.01}=2.33$. Then $P$-value $>.01$. Decision: do not reject $H_0$.
\end{itemize}
\fig{pvalue_cases}{0.85}{Upper-tailed $z$ test with $\alpha=.01$: the $P$-value is the area to the right of the calculated $z$.}

Conclusion: the $P$-value and the rejection region approach give the same decision.
Generalizing this example to other tests yields the following decision rule.

\textbf{The decision rule based on the $P$-value.} First specify $\alpha$. Then
\begin{enumerate}
  \item[(i)] if $P$-value $\le\alpha$, reject $H_0$;
  \item[(ii)] if $P$-value $>\alpha$, do not reject $H_0$.
\end{enumerate}

\textbf{The size of the $P$-value matters.} Suppose $\alpha=.05$.
\begin{itemize}
  \item Case I: $P$-value $=.0498$ $\to$ barely enough evidence, a close call.
  \item Case II: $P$-value $=.0003$ $\to$ very strong evidence, a clear decision.
\end{itemize}
Conclusion: the $P$-value tells not only the decision, but also how strong the evidence is.

\textbf{The $P$-value for a $z$ test.}
\begin{enumerate}
  \item Upper-tailed test ($H_a:\mu>\mu_0$): $P$-value $=P(Z\ge z \text{ when } \mu=\mu_0) = 1-\Phi(z)$ (area to the right of the calculated $z$).
  \item Lower-tailed test ($H_a:\mu<\mu_0$): $P$-value $=P(Z\le z \text{ when } \mu=\mu_0) = \Phi(z)$ (area to the left of the calculated $z$).
  \item Two-tailed test ($H_a:\mu\ne\mu_0$):
  \begin{align*}
  z>0:&\quad P\text{-value} = P(Z\le -z \text{ or } Z\ge z) = \Phi(-z)+1-\Phi(z) = 2[1-\Phi(z)], \tag{1}\\
  z<0:&\quad P\text{-value} = P(Z\le z \text{ or } Z\ge -z) = \Phi(z)+1-\Phi(-z) = 2[1-\Phi(-z)]. \tag{2}
  \end{align*}
  From (1) and (2), the $P$-value for the two-tailed test is
  \[
  P\text{-value} = 2[1-\Phi(|z|)],
  \]
  the sum of the two tail areas.
\end{enumerate}
\fig{pvalue_z}{0.95}{$P$-value of the $z$ test (shaded area) for the three alternatives.}


\textbf{The $P$-value for a $t$ test.}
\begin{itemize}
  \item The ideas are the same as for the $z$ test (the $P$-value is an area under the curve).
  \item Degrees of freedom matter: if the df is different, then the $t$ distributions are different, which yields different $P$-values.
  \item The decision rule is the same.
\end{itemize}

\textbf{More on interpreting the $P$-value.}
\begin{itemize}
  \item The $P$-value is \emph{not} the probability that $H_0$ is true or false.
  \item The $P$-value is \emph{not} the probability of (wrongly) rejecting $H_0$; that is $\alpha$.
  \item The $P$-value varies from sample to sample because (i) the test statistic changes from sample to sample, (ii) so it is also a random variable.
  \item If $H_0$ is false, then the $P$-value tends to be small (so we can reject $H_0$).
\end{itemize}
